\documentclass{article}

\usepackage[a4paper,landscape,margin=2cm]{geometry}
\usepackage{../../../../components/mathtex}
\usepackage{colortbl}
\usepackage{array}

\usepackage{fancyhdr}
\usepackage{ulem}
\usepackage{paracol}

\pagestyle{fancy}
\fancyhf{}
\fancyhead[L]{DL3 Math-Info}
\fancyhead[C]{Fiche d'entraînement hebdomadaire N°9}
\fancyhead[R]{2026-2027}
\fancyfoot[L]{Ewen Rodrigues de Oliveira}
\fancyfoot[R]{\thepage}


\begin{document}
\setlength{\columnsep}{2cm}

\begin{paracol}{2}

\renewcommand{\arraystretch}{1.3}
\noindent
\begin{tabular}{|p{0.6\columnwidth}|>{\centering\arraybackslash}p{0.35\columnwidth}|}
\hline
\rowcolor[RGB]{230,237,247}
\textbf{Matière} & \textbf{Exercices} \\
\hline
Intégration et probabilités & $6,7,8,9,10$\\
\hline
Groupes et actions & $1,2,3,4,5$\\
\hline
Calcul différentiel & $11,12,13$\\
\hline
\end{tabular}

\vspace{1em}

\exercise{1}{Partitions et classes de conjugaison dans $\mathcal{S}_n$}{
Soit $n \ge 1$. Étant donnée une permutation $\sigma \in \mathcal{S}_n$, on écrit $\sigma$ comme produit de cycles à supports disjoints $\sigma = \sigma_1 \dots \sigma_r$. On note $\ell_i$ la longueur du cycle $\sigma_i$ pour chaque $1 \le i \le r$.
\begin{enumerate}
    \item Expliquer pourquoi on peut supposer que $\ell_1 \le \ell_2 \le \dots \le \ell_r$. On fera cette hypothèse dans la suite de l'exercice.
    \item Soit $m$ le nombre de points fixes de $\sigma$. On associe à $\sigma$ le $(r+m)$-uplet $\mathrm{type}(\sigma) = (1, \dots, 1, \ell_1, \dots, \ell_r)$, où les $m$ premières entrées sont égales à 1. Montrer que $\mathrm{type}(\sigma)$ est une partition de $n$, c'est-à-dire un $s$-uplet $(a_1, \dots, a_s)$ (pour un certain $s \ge 1$) d'entiers strictement positifs tels que $a_1 + \dots + a_s = n$ et $a_1 \le \dots \le a_s$.
    \item Montrer que le type induit une bijection entre l'ensemble des classes de conjugaison de $\mathcal{S}_n$ et l'ensemble des partitions de $n$.
    \item On se place dans $\mathcal{S}_9$. Parmi les permutations suivantes, lesquelles sont conjuguées :
    \[
    \sigma_1 = (1 \; 6 \; 3)(2 \; 8 \; 9 \; 5), \quad \sigma_2 = (1 \; 9 \; 2)(4 \; 9 \; 8 \; 5), \quad \sigma_3 = (2 \; 6 \; 7 \; 9)(4 \; 5 \; 8) ?
    \]
    \item Combien y a-t-il de classes de conjugaisons dans $\mathcal{S}_6$ ?
\end{enumerate}
}

\exercise{2}{Sous-groupes normaux de $\mathcal{S}_4$ et du groupe alterné}{
\begin{enumerate}
    \item Faire la liste des sous-groupes normaux de $\mathcal{S}_4$ (\emph{indication : un sous-groupe normal est une réunion disjointes de classes de conjugaison}). Vérifier qu'il y en a exactement un d'ordre 4. On le note $H$.

    \item Vérifier que $H$ est contenu dans $\mathcal{A}_4$ et calculer le quotient $\mathcal{A}_4 / H$.
    \item Montrer que le morphisme obtenu en composant l'injection naturelle $\mathcal{S}_3 \to \mathcal{S}_4$ et la projection canonique $\mathcal{S}_4 \to \mathcal{S}_4 / H$ est un isomorphisme.
    \item Montrer que $D(\mathcal{A}_4) = H$.
       \switchcolumn
    \setcounter{enumi}{4}
    \item Combien y a-t-il de morphismes de $\mathcal{A}_4$ dans $\mathbb{C}^\times$ ?
\end{enumerate}
}

\exercise{3}{Étude d'un sous-groupe engendré dans $\mathcal{S}_{11}$}{
Dans le groupe symétrique $\mathcal{S}_{11}$, on considère les éléments :
\[
\sigma := (1 \; 2), \quad \tau := (1 \; 2 \; 4 \; 3)(5 \; 6 \; 7 \; 8)(10 \; 3 \; 2 \; 4 \; 8 \; 9).
\]
On pose $H := \langle \sigma, \tau \rangle$.
\begin{enumerate}
    \item Écrire $\tau$ comme produit de cycles à supports disjoints. Quel est l'ordre de $\tau$ ? Est-ce que $\tau$ appartient au groupe alterné $\mathcal{A}_{11}$ ?
    \item Reconnaître le groupe $H$. Quel est son indice dans $\mathcal{S}_{11}$ ?
    \item Montrer que, pour tout $\rho \in \mathcal{S}_{11} \setminus H$, le groupe $\mathcal{S}_{11}$ est engendré par $\sigma$, $\tau$ et $\rho$.
\end{enumerate}
}

\exercise{4}{Normalisateur et centralisateur}{
Soient $G$ un groupe et $A$ une partie de $G$.\\
Posons $N(A) = \{g \in G,\; gAg^{-1} = A\}$ et $C(A) = \{g \in G,\; \forall a \in A,\; gag^{-1} = a\}$, ce sont respectivement le \emph{normalisateur} et le \emph{centralisateur} de $A$.\\
Montrer que $N(A)$ et $C(A)$ sont des sous-groupes de $G$ et que $C(A)$ est un sous-groupe distingué de $N(A)$.
}

\exercise{5}{Groupe quotient $\mathbb{Q}/\mathbb{Z}$}{
On considère le groupe quotient $(\mathbb{Q}/\mathbb{Z}, +)$.
\begin{enumerate}
    \item Montrer que l'application $\varphi : (\mathbb{R}, +) \to (\mathbb{C}^*, \times)$ définie par $\varphi(x) = e^{2i\pi x}$ est un homomorphisme de groupes.
    \item Calculer son noyau et son image, et en déduire que $\mathbb{R}/\mathbb{Z} \simeq \mathbb{U} := \{z \in \mathbb{C} \,;\, |z| = 1\}$.

    \item On note $\phi$ la restriction de $\varphi$ au sous-groupe $\mathbb{Q}$. Montrer (par double inclusion) que l'image de $\phi$ est égale à
    \[
    \mathbb{U}_\infty = \{z \in \mathbb{C} \,;\, \text{il existe } n \in \mathbb{N}^* \text{ tel que } z^n = 1\} = \bigcup_{n \in \mathbb{N}^*} \mathbb{U}_n.
    \]
    \item En déduire que $\mathbb{Q}/\mathbb{Z} \simeq \mathbb{U}_\infty$.
\end{enumerate}
}
\switchcolumn
\newpage

\exercise{6}{Vrai/Faux}{
\begin{enumerate}
    \item Lorsque la v.a.r. $X$ ne prend qu'un nombre fini de valeurs $\{x_1, \dots, x_n\}$, l'espérance de $X$ est égale à la moyenne arithmétique de ces valeurs.
    \item Une v.a.r de carré intégrable et possédant une densité ne peut avoir une variance nulle.
    \item Une v.a.r de carré intégrable est toujours intégrable.
    \item Si la v.a.r $X$ vérifie $\mathbb{E}[\exp(X)] < \infty$ alors $X$ est intégrable.
    \item Si la v.a.r $X$ possède des moments de tous ordres et si $f$ est une fonction continue de $\mathbb{R}$ dans $\mathbb{R}$ alors $f(X)$ possède des moments de tous ordres.
    \item Pour une v.a.r $X$ telle que les espérances ci-dessous sont définies, on a toujours
    \[
    \mathbb{E}[\exp(-X)] \ge \exp(-\mathbb{E}[X]) \quad \mathbb{E}[\exp(X)] \ge \exp(\mathbb{E}[X]).
    \]
\end{enumerate}
}

\exercise{7}{Minimisation de l'erreur quadratique moyenne}{
Soit $X$ une variable aléatoire réelle de carré intégrable. Trouver la constante $m$ telle que
\[
\mathbb{E}[(X - m)^2] = \inf_{y \in \mathbb{R}} \mathbb{E}[(X - y)^2]
\]
}

\exercise{8}{Espérance d'une variable aléatoire à valeurs entières}{
Soit $X$ une variable aléatoire. Montrer que si $X$ prend ses valeurs dans $\mathbb{N}$, alors :
\[
\mathbb{E}(X) = \sum_{n=0}^{\infty} \mathbb{P}(X > n)
\]
}

\exercise{9}{Espérance d'une variable aléatoire à densité positive}{
Soit $X$ une variable aléatoire positive de densité $f$. On note $F$ sa fonction de répartition. Montrer que l'on a :
\[
\mathbb{E}(X) = \int_{0}^{\infty} (1 - F(t)) \, dt.
\]
}
\switchcolumn
\newpage

\exercise{10}{Continuité et compacts}{
On munit $\mathbb{R}^n$ de sa topologie usuelle (induite par n'importe quelle norme). Déterminer si les parties suivantes sont compactes :
\begin{enumerate}
    \item $A = [0,1] \times \{2\}$.
    \item $B = \{(x,y) \in \mathbb{R}^2 \mid x^2 + y^4 \le 3\}$. Existe-t-il une plus grande valeur de $\sin(xy) + e^{x-y}$ sur $B$ ?
    \item $C = \{(x,y) \in \mathbb{R}^2 \mid \max\{|x|, |y|\} = 2\}$.
\end{enumerate}
}

\exercise{11}{Complétude de $(\mathbb{R}^n, \|\cdot\|_\infty)$}{
Soit $E = \mathbb{R}^n$ muni de la norme $\|\cdot\|_\infty$ donnée par
\[
\forall x = (x_1, x_2, \dots, x_n) \in \mathbb{R}^n, \quad \|x\|_\infty = \max_{1 \le i \le n} |x_i|.
\]
Montrer que $(E, \|\cdot\|_\infty)$ est complet.
}

\exercise{12}{Topologie en dimension infinie}{
On considère $E = \ell^\infty$ l'espace vectoriel des suites réelles bornées. Pour tout $u = (u(k))_k \in E$ on pose :
\[ \|u\|_\infty = \sup_{k \in \mathbb{N}} |u(k)| \]
On considère deux sous-espaces vectoriels de $E$, définis par :
\[ A := \left\{ u = (u(k))_k \in E \; \middle| \; u(k) \xrightarrow[k \to +\infty]{} 0 \right\} \]
\[ B := \left\{ u = (u(k))_k \in E \; \middle| \; \exists K \in \mathbb{N}, \forall k \ge K, u(k) = 0 \right\} \]

\begin{enumerate}
    \item Montrer que $\|\cdot\|_\infty$ est une norme sur $E$.

    \item[2.] Montrer que $A$ est fermé dans $(E, \|\cdot\|_\infty)$, mais que $B$ ne l'est pas.
    
    \textit{Indication :} On pourra utiliser la suite $(u_n)_n$ d'éléments de $B$ définie par $u_n(k) = \begin{cases} 1 & \text{si } k \le n \\ 0 & \text{sinon} \end{cases}$.
    \item[3.] Montrer que $A$ est l'adhérence de $B$ dans $(E, \|\cdot\|_\infty)$.
\end{enumerate}
}

\end{paracol}
\end{document}